% 余白の設定
\setlength{\textwidth}{\fullwidth}
\setlength{\textheight}{40\baselineskip}
\addtolength{\textheight}{\topskip}
\setlength{\voffset}{-0.2in}
\setlength{\topmargin}{0pt}
\setlength{\headheight}{0pt}
\setlength{\headsep}{0pt}

% パッケージ
\usepackage[dvipdfmx]{hyperref,graphicx}		% 画像
\hypersetup{
	colorlinks=true, % リンクに色をつけない設定
	bookmarks=true, % 以下ブックマークに関する設定
	bookmarksnumbered=true,
	pdfborder={0 0 0},
	bookmarkstype=toc
}

\usepackage{ascmac}
\usepackage{amsmath, amssymb}		% ギリシャ文字 % 数式色々
% \usepackage{amsfonts} % \mathbb{A}
\usepackage{mathtools}
\usepackage{bm}				% 数式　\bm{a} aベクトル
\usepackage{comment}			% コメント
\usepackage{siunitx}			% SI単位
\usepackage{framed}			% 枠組み
\usepackage{braket}     % ブラケット記法
\usepackage[square,sort,comma,numbers]{natbib}


\newenvironment{Theorem}{\begin{screen}\begin{theorem}}{\end{theorem}\end{screen}}
\newenvironment{Prop}{\begin{screen}\begin{prop}}{\end{prop}\end{screen}}
\newenvironment{Lemma}{\begin{screen}\begin{lemma}}{\end{lemma}\end{screen}}
\newenvironment{Corr}{\begin{screen}\begin{corr}}{\end{corr}\end{screen}}
\newenvironment{Reidai}{\begin{screen}\begin{reidai}}{\end{reidai}\end{screen}}
\newenvironment{Def}{\begin{shadebox}\begin{DEF}}{\end{DEF}\end{shadebox}}
\newenvironment{Prob}{\begin{screen}\begin{mondai}}{\end{mondai}\end{screen}}
\newenvironment{Memo}{\begin{boxnote}\begin{memo}}{\end{memo}\end{boxnote}}

\newtheorem{prop}{性質}
\newtheorem{lemma}{補題}
\newtheorem{corr}{系}
\newtheorem{theorem}{定理}
\newtheorem{DEF}{定義}
\newtheorem{remark}{注意}
\newtheorem{reidai}{{例題}}
\newtheorem{mondai}{{問題}}
\newtheorem{memo}{{メモ}}

\newcommand{\nequiv}{\equiv\hspace{-3mm}/\; }
\newcommand{\proof}{\noindent 証明:\ }
\newcommand{\sol}{\noindent 解:\ }
\newcommand{\qed}{\qquad $\blacksquare$}

\newcommand{\K}{\mathbb{K}}
\newcommand{\R}{\mathbb{R}}
\newcommand{\C}{\mathbb{C}}

\newcommand{\norm}[1]{\|#1\|}

\newcommand{\BF}[1]{{\bf#1}}



\begin{comment}
　複数行に渡るコメントの書き方
　\usepackage{comment}が必要。
\end{comment}

% section前に改ページ
\makeatletter
\def\section{\newpage\@startsection {section}{1}{\z@}{-3.5ex plus -1ex minus -.2ex}{2.3 ex plus .2ex}{\Large\bf}}
\makeatother

% 数式番号をいい感じに
\makeatletter
% \renewcommand{\theequation}{\arabic{chapter}-\arabic{section}-\arabic{equation}}
	\renewcommand{\theequation}{\arabic{section}-\arabic{equation}}
  \@addtoreset{equation}{section}
\makeatother
